Existing frameworks like \textbf{AutoGPT} \cite{siggrav2023autogpt} and \textbf{LangGraph} \cite{langgraph2023} focus on iterative refinement or static task graphs. However, they lack an adversarial validation layer and the ability to structurally evolve their own logic over generations without human intervention. Recent advancements in \textbf{Agentic Workflows} \cite{ng2024agentic} emphasize design patterns, while \textbf{RL-PPO} based methods optimize for narrow reward functions \cite{ouyang2022training}. In contrast, our protocol optimizes for systemic stability and logical fidelity.

While \textbf{Genetic Algorithms} \cite{goldberg1989genetic} and \textbf{Artificial Immune Systems} (AIS) \cite{de2002artificial, dasgupta2011artificial} have been long established, they typically rely on random bit-wise mutations or fixed antigen-antibody matching rules. Stephanie Forrest's pioneering work on self-nonself discrimination \cite{forrest1994self} provides the conceptual basis for our Digital Thymus, yet our approach introduces semantic understanding through LLM-driven ``T-Cell'' audits, bridging the gap between symbolic AI and connectionist models.

\subsection{Safe AI Evolution \& Constitutional AI}
\textbf{Constitutional AI} \cite{bai2022constitutional} uses fixed recursive self-improvement principles. The Dialectical Protocol allows the ``Genetic Priors'' (The Constitution) to evolve based on performance metadata in a sandbox, governed by \textbf{Merkle-root validation} to prevent catastrophic drift. This concept is further supported by recent work on LLMs as optimizers \cite{yang2023llm_optimizer}.
